% ============================================================================
%  chapter7.tex
%  Inter-Process Communication: Pipes and the Art of Wiring Processes
% ============================================================================

\chapter{Inter-Process Communication: Pipes}

\section{Why Processes Cannot Just Talk}

Chapter~5 gave each process a \emph{private} address space and told you
the kernel forbids one process from poking another's memory. That
isolation is the price of safety. But isolation would be useless if
processes could not cooperate. A compiler must feed its output to an
assembler; a web server must hand a request to a worker; a database
must talk to its replicas. The general problem is called
\emph{inter-process communication} (IPC): \emph{how do two independent
processes exchange data when they cannot touch each other's memory?}

\begin{verbatim}
   What processes CANNOT do:            What the OS provides instead:

   Process A reads Process B's heap.    channels that cross the kernel:
        A -X-> B                         pipes, files, shared memory,
                                         sockets, signals (ch.6)

   Both point at "the same" variable.   The OS copies bytes across the
   No shared addresses exist.            boundary; the kernel does the
                                         moving.
\end{verbatim}

There are five classic IPC mechanisms, each a different trade-off, and
we will meet them all by the end of the book. This chapter builds the
first and most elegant one---the \emph{pipe}---because it is the
simplest possible answer and because it is the mechanism hiding behind
every command line you have ever piped.

\begin{osconcept}
IPC is an operating system problem for a reason that flows directly
from Chapter~5: processes are isolated by the kernel, so only the
kernel can connect them. A channel must be a kernel object---something
both processes name through descriptors---because the kernel is the
only party both can trust and the only party with the power to move
bytes between private address spaces. IPC is the deliberate exception
to the isolation rule, and the OS decides exactly which exceptions are
allowed.
\end{osconcept}

\section{What a Pipe Is}

A \emph{pipe} is a unidirectional byte stream connecting two endpoints.
You create one with \code{pipe()}, which gives you a pair of
descriptors: \code{pfd[0]} is the \emph{read end}, \code{pfd[1]} is
the \emph{write end}. Bytes written to \code{pfd[1]} appear, in order,
at \code{pfd[0]}.

\begin{lstlisting}[caption={A pipe is two descriptors: read end and write end.}]
#include <stdio.h>
#include <stdlib.h>
#include <unistd.h>

int main(void) {
    int pfd[2];

    if (pipe(pfd) < 0) {
        perror("pipe");
        exit(1);
    }

    printf("read end  = fd %d\n", pfd[0]);
    printf("write end = fd %d\n", pfd[1]);

    /* Write two bytes, then close the write end to signal EOF. */
    write(pfd[1], "hi", 2);
    close(pfd[1]);

    {
        char buf[8];
        ssize_t got;

        got = read(pfd[0], buf, sizeof(buf));
        if (got < 0) {
            perror("read");
            exit(1);
        }
        printf("read %zd bytes: %.*s\n", got, (int) got, buf);
    }

    close(pfd[0]);
    return 0;
}
\end{lstlisting}

\begin{verbatim}
   The pipe is not a file. It is a KERNEL buffer with two doors:

        write end (fd[1])                read end (fd[0])
             |                                ^
             v                                |
        +--------------------------------------------+
        |  [ 'h' ][ 'i' ][ ][ ][ ][ ]  ...          |
        |         the kernel's pipe buffer           |
        +--------------------------------------------+

   - "Byte stream" means there are NO message boundaries: whatever the
     writer writes is a continuous sequence of bytes to the reader. If
     you write "ab" and "cd", the reader may see "abcd", "abc"+"d", or
     "a"+"bcd" depending on timing. Boundaries are NOT preserved.
   - The pipe is UNIDIRECTIONAL: bytes flow only from write end to
     read end. For two-way conversation you create two pipes.
   - The bytes live in the KERNEL. The processes never share memory;
     they only pass bytes across the syscall boundary.
\end{verbatim}

Two properties of this picture matter enormously. First, the pipe
synchronizes the two sides \emph{automatically}: the kernel makes one
side wait for the other. Second, the pipe has a finite capacity---on
Linux, around 64 kilobytes---and that capacity, not the file system, is
what makes blocking real.

\begin{osconcept}
A pipe is a kernel object: kernel memory for the buffer, kernel
bookkeeping for the two endpoints, and kernel rules about when a
process may proceed. Your program never sees the buffer; it sees two
integers. This is the file-descriptor model of Chapter~5 applied to
live data flow. Only the kernel can own the buffer, because only the
kernel can mediate between two isolated address spaces without copying
through either one's memory.
\end{osconcept}

\section{The Blocking Contract}

Pipes are \emph{blocking}, and that is their superpower. The kernel
makes the two ends wait for each other:

\begin{verbatim}
   EMPTY pipe:  read() blocks until bytes arrive.

        write end                 read end
             |                       v  (waiting...)
        +------------------------------------+
        |                                    |   <-- nothing to read
        +------------------------------------+
             (writer writes)
        +------------------------------------+
        | [ 'h' ][ 'i' ][ ][ ][ ][ ]        |
        +------------------------------------+
        read() wakes up and returns the bytes.

   FULL pipe:   write() blocks until there is room.

        +------------------------------------+
        | [X][X][X][X][X][X][X][X]           |   <-- no room
        +------------------------------------+
             (reader reads)
        +------------------------------------+
        | [X][X][X][X][ ][ ][ ][ ]          |
        +------------------------------------+
        write() wakes up and deposits its bytes.
\end{verbatim}

This is worth slowing down for, because it is the whole point. The
kernel is acting as a \emph{scheduler}: the reader does not busy-wait
spinning the CPU; it goes to sleep and the kernel wakes it exactly when
data is available. The writer sleeps when the buffer is full and wakes
when the reader frees space. Neither process polls, and neither needs
to know the other exists.

\begin{verbatim}
   The producer/consumer dance:

   producer:                    consumer:
     write(fd1, data, n);         got = read(fd0, buf, n);
       |                              ^
       |  if the pipe is full,        |  if the pipe is empty,
       |  the kernel parks us         |  the kernel parks us
       |  and wakes us later          |  and wakes us later
       v                              |
   bytes cross the kernel buffer, paced by whichever side is slower.
\end{verbatim}

The blocking contract is also where the classic \emph{deadlock} lives:
if two processes each write into a pipe the other is supposed to read,
and both pipes are full, each process sits blocked forever, each
waiting for the other to make room. The kernel cannot rescue you; it
only enforces the waiting.

\begin{osconcept}
Blocking I/O is the kernel's oldest concurrency solution. By parking a
process that cannot make progress, the OS wastes no CPU and needs no
coordination from the program; by waking it the instant it can proceed,
the OS delivers an efficient rendezvous. Deciding when to sleep and
when to wake---for every pipe, on every byte---is the scheduler's
problem, and it is exactly why a pipe is an OS service rather than a
library function.
\end{osconcept}

\section{fork + pipe: Wiring Two Processes Together}

A pipe inside one process is a curiosity. The pipe becomes IPC the
moment you \code{fork}: the child inherits a \emph{copy} of the
descriptor table, so both processes hold the same two ends. One process
keeps the write end, the other keeps the read end, and bytes flow
between two independent address spaces.

\begin{lstlisting}[caption={Parent writes, child reads: the canonical IPC dance.}]
#include <stdio.h>
#include <stdlib.h>
#include <unistd.h>
#include <sys/types.h>
#include <sys/wait.h>

int main(void) {
    int pfd[2];
    pid_t pid;

    if (pipe(pfd) < 0) {
        perror("pipe");
        exit(1);
    }

    pid = fork();
    if (pid < 0) {
        perror("fork");
        exit(1);
    }

    if (pid == 0) {
        /* CHILD: keep the read end, drop the write end. */
        char buf[64];
        ssize_t got;

        close(pfd[1]);

        got = read(pfd[0], buf, sizeof(buf));
        if (got < 0) {
            perror("read");
            exit(1);
        }
        printf("child read %zd bytes: %.*s\n", got, (int) got, buf);

        close(pfd[0]);
        exit(0);
    }

    /* PARENT: keep the write end, drop the read end. */
    close(pfd[0]);

    write(pfd[1], "hello from parent", 17);

    close(pfd[1]);            /* closing the last write end sends EOF */

    waitpid(pid, NULL, 0);
    return 0;
}
\end{lstlisting}

\begin{verbatim}
   The state after fork, before each side closes its unused end:

     parent                          child
   +-------------+                 +-------------+
   | fd table    |                 | fd table    |
   | 0 stdin     |                 | 0 stdin     |
   | 1 stdout    |                 | 1 stdout    |
   | 2 stderr    |                 | 2 stderr    |
   | 3 pfd[0] --+----------------->| 3 pfd[0]    |   read end
   | 4 pfd[1] --+----------------->| 4 pfd[1]    |   write end
   +-------------+                 +-------------+

   The kernel pipe object is shared: it is reference-counted, like the
   open file table. BOTH processes name it, so EITHER can read or write.

   After each side closes its unused end, the roles are clean:

     parent: [write end] =============> [read end]: child
                 \ byte stream flows through the kernel /
\end{verbatim}

The discipline here is absolute and it is the number-one pipe bug in
all of systems programming: \emph{close the end you do not use, in
every process}. If the parent keeps the read end open while the child
waits for EOF, the EOF never arrives (the kernel still counts the
parent's copy as a reader), and the child hangs forever.

\begin{ruleofthumb}[Close both ends you don't use, in both processes]
After \code{fork}, each process must close the pipe end it will not
use. Leaving a spare copy of the read end open is the classic way to
make a reader wait forever for an EOF that can never come. Draw the
descriptor table for every process in your program; there are never
extra descriptors on purpose.
\end{ruleofthumb}

\section{The EOF Rule}

When does \code{read} on a pipe return 0, signaling end of stream? The
rule is precise: \emph{when every copy of the write end is closed}.
EOF is not ``the writer has nothing more to say right now''; it is
``no writer exists anymore, anywhere.''

\begin{verbatim}
   EOF rule: read() returns 0 only when NO write end remains open.

   writer open                            all writers gone
        |                                      |
        v                                      v
   read() returns bytes...            read() returns 0  ("end of stream")

   This is why the child in the previous program could finish: the
   parent wrote its bytes AND closed pfd[1]. If the parent had left
   pfd[1] open, the child's read() would sit forever, even after the
   last byte, because a writer still existed.
\end{verbatim}

The EOF rule is what makes pipes compose into pipelines: a downstream
process knows the input is finished because the upstream process
\emph{closed its end}. The closing, not any special marker byte, is
the signal.

\begin{osconcept}
The EOF rule is a kernel decision about reference counting. The kernel
tracks how many open descriptors name each pipe endpoint, and it is the
one authority that knows when the count reaches zero. Your program
cannot ``force'' an EOF; it can only close descriptors and let the
kernel draw the conclusion. This is descriptor-table bookkeeping
exposed as a data-flow event, and it is what lets processes signal
completion without any extra protocol.
\end{osconcept}

\section{Pipelines: cmd1 | cmd2}

With the EOF rule and \code{dup2} from Chapter~5 in hand, we can build
the thing that made Unix famous: the \emph{pipeline}. ``Run the first
program with its stdout connected to the second program's stdin.''
The recipe has exactly four moves: create a pipe, fork the first
command with its stdout re-pointed at the write end, fork the second
with its stdin re-pointed at the read end, and have the parent close
both pipe ends and wait.

\begin{lstlisting}[caption={A two-stage pipeline, end to end.}]
#include <stdio.h>
#include <stdlib.h>
#include <unistd.h>
#include <sys/types.h>
#include <sys/wait.h>

/* Run "prog1..." with its stdout piped into "prog2...". */
int run_pipeline(char *const argv1[], char *const argv2[]) {
    int pfd[2];
    pid_t p1;
    pid_t p2;

    if (pipe(pfd) < 0) {
        perror("pipe");
        return -1;
    }

    p1 = fork();
    if (p1 < 0) {
        perror("fork");
        return -1;
    }
    if (p1 == 0) {
        /* First command: stdout must name the pipe's write end. */
        close(pfd[0]);
        if (dup2(pfd[1], 1) < 0) {
            perror("dup2");
            exit(1);
        }
        close(pfd[1]);
        execvp(argv1[0], argv1);
        perror("execvp");
        exit(127);
    }

    p2 = fork();
    if (p2 < 0) {
        perror("fork");
        return -1;
    }
    if (p2 == 0) {
        /* Second command: stdin must name the pipe's read end. */
        close(pfd[1]);
        if (dup2(pfd[0], 0) < 0) {
            perror("dup2");
            exit(1);
        }
        close(pfd[0]);
        execvp(argv2[0], argv2);
        perror("execvp");
        exit(127);
    }

    /* The parent is only the referee: it needs neither pipe end. */
    close(pfd[0]);
    close(pfd[1]);

    waitpid(p1, NULL, 0);
    waitpid(p2, NULL, 0);
    return 0;
}

int main(int argc, char *argv[]) {
    char *cmd1[3];
    char *cmd2[3];

    if (argc != 5) {
        fprintf(stderr, "usage: %s prog1 a1 prog2 a2\n", argv[0]);
        return 1;
    }
    cmd1[0] = argv[1];
    cmd1[1] = argv[2];
    cmd1[2] = NULL;
    cmd2[0] = argv[3];
    cmd2[1] = argv[4];
    cmd2[2] = NULL;

    return run_pipeline(cmd1, cmd2);
}
\end{lstlisting}

\begin{verbatim}
   The wiring, drawn process by process:

     p1 (child 1)                  p2 (child 2)
     fd 1 --> pipe write end       fd 0 --> pipe read end
           \                        /
            \  +----------------+ /
             \ | kernel buffer  |/
              >| pfd[1] -> pfd[0]|
              <|                |<   bytes flow p1 -> p2
               +----------------+

     parent: closes both ends; just waitpid()s both children.
     dup2() copies the pipe descriptor ONTO descriptors 1 and 0, so
     exec'd programs write to stdout / read from stdin without knowing
     a pipe exists -- Chapter 5's redirection, exactly as the shell
     does it.
\end{verbatim}

The magic ingredient is that \code{exec} preserves open descriptors:
after \code{dup2} and \code{close}, the child \code{exec}s a totally
new program whose \code{stdout} is the pipe, with no knowledge of the
plumbing. That is why every Unix tool composes with every other tool:
each reads \code{stdin} and writes \code{stdout}, and the shell---or
your program---rewires the descriptors.

\begin{verbatim}
   The full lifecycle of "cat f.txt | sort":

     1. shell pipe()s -> pfd = [3, 4]
     2. shell forks child A: dup2(4, 1), close(3), close(4), exec cat
        (cat's stdout is now the pipe)
     3. shell forks child B: dup2(3, 0), close(3), close(4), exec sort
        (sort's stdin is now the pipe)
     4. shell closes 3 and 4, then waits for both
     5. cat writes f.txt into the pipe; sort reads it, sorts, prints;
        when cat exits and its stdout closes, the last write end dies,
        sort's read() hits EOF, sort prints and exits, shell reaps both.
\end{verbatim}

\begin{osconcept}
Pipelines work because the OS treats descriptors as \emph{rewireable
connections} whose ownership can be transferred between processes and
onto the standard slots. The kernel's descriptor-copying primitive
(\code{dup2}) plus the EOF rule are the entire protocol---no message
headers, no handshaking. Teaching processes to ``just write to fd 1''
and letting the kernel route it is a design decision so good that it
became the model for every later IPC mechanism.
\end{osconcept}

\section{SIGPIPE: When the Channel Dies}

What happens if the reader is gone? Suppose the second command in a
pipeline exits early, closing the read end, while the first command
keeps writing. The kernel detects that a write has no readers and does
two things: it sends \code{SIGPIPE} to the writer, and if the signal
is ignored or handled, the \code{write} returns \code{-1} with
\code{errno = EPIPE}.

\begin{verbatim}
   reader closes its read end, exits
        |
        v
   writer calls write(fd1, ...)
        |
        +--> kernel finds no readers
        |
        +--> raises SIGPIPE on the writer
        |      default action: terminate the writer (like `yes | head`)
        |
        +--> if SIGPIPE ignored/handled:
             write() returns -1, errno = EPIPE

   This is why `yes | head -1` terminates `yes` promptly: `head` closes
   the pipe after one line, `yes` writes, the kernel SIGPIPEs it.
\end{verbatim}

For a well-behaved server or filter, the sensible choice is to ignore
\code{SIGPIPE} and treat \code{EPIPE} as ``the peer went away''---an
ordinary, expected condition, not a crash.

\begin{lstlisting}[caption={Handling a dead reader gracefully.}]
#include <stdio.h>
#include <stdlib.h>
#include <signal.h>
#include <unistd.h>
#include <errno.h>

int main(void) {
    struct sigaction sa;
    char msg[16];
    int i;

    /* The reader may vanish; treat that as EPIPE, not death. */
    sa.sa_handler = SIG_IGN;
    sigemptyset(&sa.sa_mask);
    sa.sa_flags = 0;
    sigaction(SIGPIPE, &sa, NULL);

    /* Write a line repeatedly until the reader goes away. */
    for (i = 0; i < 100000; i++) {
        ssize_t r;

        snprintf(msg, sizeof(msg), "line %d\n", i);
        r = write(1, msg, strlen(msg));
        if (r < 0) {
            if (errno == EPIPE) {
                printf("reader is gone; stopping politely\n");
                break;
            }
            perror("write");
            break;
        }
    }
    return 0;
}
\end{lstlisting}

\begin{osconcept}
\code{SIGPIPE} is the kernel's way of converting a dead channel into
an interruptible event. The OS cannot know whether a missing reader is
a bug or an expected teardown, so it gives the writer the choice:
die by default, or catch the signal and inspect \code{EPIPE}. Choosing
the default is the reason so many quick pipelines terminate cleanly;
choosing the handler is what production services do so that a client
disconnecting is not a crash.
\end{osconcept}

\section{Named Pipes (FIFOs)}

An unnamed pipe lives and dies with the processes that created it. A
\emph{named pipe}, or \emph{FIFO}, is a pipe with a name in the file
system. Any process that can open the path can connect to it, so two
unrelated processes---neither forking the other---can talk through a
FIFO. You create one with \code{mkfifo} (the program or the syscall);
you open it like a file; the data still flows through a kernel buffer,
not through the file system.

\begin{lstlisting}[caption={Creating and opening a FIFO.}]
#include <stdio.h>
#include <stdlib.h>
#include <sys/stat.h>
#include <fcntl.h>
#include <unistd.h>

int main(int argc, char *argv[]) {
    int fd;

    if (argc != 2) {
        fprintf(stderr, "usage: %s fifo_name\n", argv[0]);
        return 1;
    }

    /* Create the FIFO if it does not exist yet. */
    if (mkfifo(argv[1], 0644) < 0) {
        perror("mkfifo");   /* fine if it already exists (EEXIST) */
    }

    printf("opening %s for writing (blocks until a reader opens it)...\n",
           argv[1]);
    fd = open(argv[1], O_WRONLY);
    if (fd < 0) {
        perror("open");
        return 1;
    }

    write(fd, "hello through a named pipe", 26);
    close(fd);
    return 0;
}
\end{lstlisting}

\begin{verbatim}
   FIFO = a pipe with an address in the file namespace:

      $ mkfifo /tmp/chat
      $ ./writer /tmp/chat &     (blocks: no reader yet)
      $ cat /tmp/chat            (reader opens the other end)
      hello through a named pipe

   - The FIFO is a directory entry, but it holds NO data of its own.
     Bytes go through the kernel pipe buffer; the path is just the
     meeting point.
   - open() BLOCKS until both a reader and a writer have opened the
     FIFO. That rendezvous is the FIFO's version of pipe() + fork().
   - When all readers or writers close, the pipe state resets: a later
     open() gets a fresh buffer, so data written earlier is gone.
\end{verbatim}

\begin{osconcept}
A FIFO is the file system serving as a \emph{namespace} for IPC. The
kernel already owns the directory tree, permissions, and lookup; giving
a pipe a path reuses all of that machinery so that unrelated processes
can rendezvous by name. Deciding that the path is only an address---not
a store of data---is an OS design choice that keeps FIFOs cheap and
their semantics identical to anonymous pipes.
\end{osconcept}

\section{The IPC Menu}

Pipes are one of five classic channels. The others appear later in the
book; a map now will save you confusion when we meet them.

\begin{verbatim}
   mechanism        data model          who can talk        we meet it
   --------------   ----------------    ----------------    ----------
   pipe / FIFO      byte stream,        parent-child or     this chapter
                    kernel-buffered     any named pair
   signal           tiny events         any processes       chapter 5
                    (an integer)
   shared memory    raw bytes, no       processes that      chapter 11
   (mmap)           kernel copy         agree on a region
   message queue    typed messages,     any processes       (surveyed)
                    kernel-buffered     with permissions
   socket           byte stream or      any processes,      chapter 16
                    datagrams           even across hosts

   The trade-off in one line: pipes and sockets BUFFER and COPY (safe,
   automatic), while shared memory does NOT (fast, dangerous). You
   choose based on whether you want the kernel to babysit your data.
\end{verbatim}

\begin{verbatim}
   The rule of thumb for this book's projects:

     one writer, one reader, stream data  -> pipe
     unrelated processes, named channel   -> FIFO
     lots of small structured messages    -> message queue
     fast shared state, both are careful  -> shared memory (ch.12)
     any machine, any distance            -> sockets (ch.17)
\end{verbatim}

\begin{researchnote}[Pipes are communication as mathematics]
The pipe is not just a Unix convenience; it is the direct implementation
of a celebrated idea. In 1978, C.~A.~R.~Hoare published
``Communicating Sequential Processes,'' arguing that the cleanest way
to think about concurrency is as \emph{processes that exchange data by
message passing}, with no shared state at all. A pipe is exactly that:
two processes, one channel, bytes in order.

\begin{verbatim}
   Hoare's thesis, condensed to one line:

       "Programs are built from sequential processes that interact
        solely by input and output; shared memory is unnecessary and
        harmful."

   This is the philosophy the pipe implements: the writer and reader
   never share a variable, never lock anything, and are correct purely
   because the channel delivers bytes in order. Hoare's later book of
   the same name turned this into a formal language for reasoning
   about concurrency.
\end{verbatim}

Why it matters to you: the same idea drives most real servers
(Chapter~17) and the concurrency chapter (Chapter~11). When a design
uses pipes or sockets to connect worker processes, it is inheriting a
four-decade-old research result about how to make concurrency
\emph{provable} instead of merely probable.

\paperref{C.~A.~R.~Hoare, ``Communicating Sequential Processes,''
Communications of the ACM 21(8), 1978}{https://www.cs.cmu.edu/~crary/819-f09/Hoare78.pdf}
\end{researchnote}

\section{Practical: A Two-Stage Filter Shell}

Let us wire the chapter together into a small program that behaves like
a real shell pipeline: read a command line like
\code{prog1 arg prog2 arg}, split it, and run it through a pipe---the
exact four-move recipe, with every error checked and every unused
descriptor closed.

\begin{lstlisting}[caption={piperun.c: parse two commands and pipe them together.}]
#include <stdio.h>
#include <stdlib.h>
#include <string.h>
#include <unistd.h>
#include <sys/types.h>
#include <sys/wait.h>

#define MAX_ARG 32

/* Split "cmd arg1 arg2 | cmd arg1" into two NULL-terminated argv arrays. */
int parse_pipeline(char *line, char ***left, char ***right) {
    char *pipe_pos;
    char **a;
    char *tok;
    int n;

    pipe_pos = strchr(line, '|');
    if (pipe_pos == NULL) {
        return -1;      /* no pipe found */
    }
    *pipe_pos = '\0';

    a = malloc(MAX_ARG * sizeof(char *));
    if (a == NULL) {
        return -1;
    }
    n = 0;
    tok = strtok(line, " \t\n");
    while (tok != NULL && n < MAX_ARG - 1) {
        a[n] = tok;
        n = n + 1;
        tok = strtok(NULL, " \t\n");
    }
    a[n] = NULL;
    *left = a;

    a = malloc(MAX_ARG * sizeof(char *));
    if (a == NULL) {
        return -1;
    }
    n = 0;
    tok = strtok(pipe_pos + 1, " \t\n");
    while (tok != NULL && n < MAX_ARG - 1) {
        a[n] = tok;
        n = n + 1;
        tok = strtok(NULL, " \t\n");
    }
    a[n] = NULL;
    *right = a;

    return 0;
}

/* Run the pipeline built by parse_pipeline. */
int run_pipeline(char *const left[], char *const right[]) {
    int pfd[2];
    pid_t p1;
    pid_t p2;

    if (pipe(pfd) < 0) {
        perror("pipe");
        return -1;
    }

    p1 = fork();
    if (p1 < 0) {
        perror("fork");
        return -1;
    }
    if (p1 == 0) {
        close(pfd[0]);
        dup2(pfd[1], 1);
        close(pfd[1]);
        execvp(left[0], left);
        perror("execvp");
        exit(127);
    }

    p2 = fork();
    if (p2 < 0) {
        perror("fork");
        return -1;
    }
    if (p2 == 0) {
        close(pfd[1]);
        dup2(pfd[0], 0);
        close(pfd[0]);
        execvp(right[0], right);
        perror("execvp");
        exit(127);
    }

    close(pfd[0]);
    close(pfd[1]);

    waitpid(p1, NULL, 0);
    waitpid(p2, NULL, 0);
    return 0;
}

int main(void) {
    char line[512];
    char **left;
    char **right;

    printf("piperun> ");
    fflush(stdout);
    while (fgets(line, sizeof(line), stdin) != NULL) {
        if (parse_pipeline(line, &left, &right) == 0) {
            run_pipeline(left, right);
            free(left);
            free(right);
        } else {
            printf("need two commands separated by a pipe\n");
        }
        printf("piperun> ");
        fflush(stdout);
    }
    return 0;
}
\end{lstlisting}

\begin{verbatim}
   Try it:

     $ ./piperun
     piperun> echo hello world | tr a-z A-Z
     HELLO WORLD
     piperun> seq 10 | tail -3
     8
     9
     10
     piperun> ^D

   Every line you typed exercised the whole chapter: a kernel buffer,
   two forks, two dup2 rewirings, four closes, and two waits. That is
   a shell's pipeline in miniature.
\end{verbatim}

\section{Structures in Bytes}

A pipe is a kernel object, so its byte structure lives in the kernel;
in your process a pipe is just two small integers. The object they
name looks like this (Linux):

\begin{verbatim}
   the pipe: a ring buffer of 16 pages = 64 KB.

   +---------+---------+---------+-- ... --+---------+
   | page 0  | page 1  | page 2  |          | page 15 |
   +---------+---------+---------+-- ... --+---------+
        ^ read position               ^ write position
        (byte offsets into the ring)

   the two file objects name the same pipe:

   file object (read end)  \   +---------------------+
   file object (write end) /-->| pipe: 16 pages      |
                              | read pos            |
                              | write pos           |
                              | nreaders, nwriters  |  <- the EOF rule
                              +---------------------+
\end{verbatim}

Writes land at the write position, wrap around the ring, and advance
it; reads do the same at the read position. The reference counts are
the byte-level truth behind the EOF rule: \code{read()} returns 0 not
because a marker byte exists, but because the \code{nwriters} counter
reached zero. The ring is also why blocking is cheap---a writer with no
room sleeps until a reader advances the read position.

\section{Summary}

\begin{verbatim}
   The pipe model, complete:

   pipe()      : one call, two descriptors, one kernel buffer
   fork()      : duplicates the descriptors into both processes
   close()     : each side drops the end it does not use
   blocking    : the kernel parks readers on empty, writers on full
   EOF         : read() returns 0 when the last write end closes
   dup2()      : rewires a descriptor, building pipelines
   SIGPIPE     : the kernel tells a writer its reader is gone
   FIFO        : the same channel, given a name in the file system

   Rules that keep pipe programs correct:
     1. close the unused end in EVERY process (the EOF killer);
     2. the parent (or middleman) closes BOTH ends;
     3. remember pipes are byte streams: no message boundaries;
     4. expect blocking: it is the feature, not a bug;
     5. handle SIGPIPE (or EPIPE) when a peer may disconnect.
\end{verbatim}

\begin{osconcept}
This chapter was a complete answer to ``how do isolated processes
talk?'' The OS answers: through channels it alone owns---buffers that
synchronize, reference counts that produce EOF, signals that report a
dead peer, and names that let strangers rendezvous. Every later IPC
mechanism is a variation on these same ideas, and the discipline
(close your ends, check your returns) never changes.
\end{osconcept}

The pipes we just built are \emph{synchronous}: the writer is paced by
the reader. In the chapters ahead we make communication more powerful
and more dangerous: threads sharing one address space (Chapter~11),
shared memory with no kernel copy (Chapter~12), and finally sockets
that cross machine boundaries (Chapter~17).

